% Generated draft: requires mathematical and source review.
\documentclass[11pt]{article}
\usepackage[margin=25mm]{geometry}
\usepackage{fontspec}
\setmainfont{Latin Modern Roman}
\usepackage{amsmath,amssymb,mathtools,unicode-math}
\setmathfont{Latin Modern Math}
\usepackage{enumitem,needspace,xurl,hyperref}
\setlength{\parindent}{0pt}
\setlength{\parskip}{5pt}
\setlength{\emergencystretch}{3em}
\newcommand{\N}{\mathbb{N}}
\newcommand{\Z}{\mathbb{Z}}
\newcommand{\Q}{\mathbb{Q}}
\newcommand{\R}{\mathbb{R}}
\newcommand{\C}{\mathbb{C}}
\newcommand{\F}{\mathbb{F}}
\newcommand{\A}{\mathbb{A}}
\newcommand{\PP}{\mathbb{P}}
\newcommand{\E}{\mathbb{E}}
\newcommand{\eps}{\varepsilon}
\newcommand{\dd}{\,\mathrm{d}}
\begin{document}
\section*{UnsolvedMath ID: 30004937}
\textbf{Source problem code:} catalogue entry
\subsection{Definitions and mathematical statement}
For integers \(1\leq k\leq n\), let \([n]=\{1,\ldots,n\}\), and let \(\operatorname{Gr}(k,n)\) be the Grassmannian of \(k\)-dimensional real linear subspaces of \(\R^n\). Its totally nonnegative part is
\[
\operatorname{Gr}_{\geq0}(k,n)=\left\{\operatorname{rowspan}(M): M\in\R^{k\times n},\ \operatorname{rank}(M)=k,\ \Delta_I(M)\geq0\text{ for every }I\in\binom{[n]}k\right\},
\]
where \(\Delta_I(M)\) denotes the maximal minor with column indices \(I\) in increasing order. Use the convention that \(\operatorname{Gr}_{\geq0}(0,n)\) consists of the zero subspace. These bounds describe when the adjacent Grassmannians are defined; the supplied source does not specify the parameter range intended for its conjecture.

For \(X\in\operatorname{Gr}_{\geq0}(k,n)\), define its Pl\"ucker support by
\[
\mathcal B(X)=\left\{I\in\binom{[n]}k:\Delta_I(M)>0\right\},
\]
using any representative \(M\) with nonnegative maximal minors. A positroid cell is a nonempty stratum on which \(\mathcal B(X)\) is fixed. The source denotes such cells by \(\Pi_f^{>0}\), using permutation indices \(f\), and describes a shifted index by
\[
f^{\downarrow}(i)=f(i)-1.
\]
The domain, periodicity, and indexing conventions for these permutations are not supplied. Thus this formula alone does not specify a valid target index for every source cell.

\textbf{Source conjecture, with unresolved specifications.} The source proposes a shift map, also called T-duality, between certain subsets
\[
S:A_{k,n}\longrightarrow B_{k-1,n},\qquad
A_{k,n}\subseteq\operatorname{Gr}_{\geq0}(k,n),\qquad
B_{k-1,n}\subseteq\operatorname{Gr}_{\geq0}(k-1,n),
\]
and asks whether this intended map is a bijection that sends every applicable source positroid cell onto its shifted target cell:
\[
S\bigl(\Pi_f^{>0}\bigr)=\Pi_{f^{\downarrow}}^{>0}.
\]
This equality records the source's intended whole-cell correspondence wherever the relevant cells belong to the intended subsets. The record does not identify those subsets, specify which cells are applicable, define \(S\) on Grassmannian points, or give its boundary behavior. Consequently this is a qualified statement of the proposed target, not a fully specified conjecture about a defined map. In particular, the dashed arrow in the source does not supply a definition or establish rationality.

The source also reports a conjectured topological equivalence between spaces of planar Ising networks and electrical networks, and says that the shift map is straightforward to define on critical planar Ising and electrical networks. These are source-specific background qualifications: no network spaces, topology, criticality conditions, or restricted map formula are supplied, so they do not resolve the missing Grassmannian formulation.
\subsection{Short English statement}
The proposed shift map should biject certain subsets of adjacent real totally nonnegative Grassmannians and carry each applicable positroid cell to the cell obtained by lowering its permutation values by one. The supplied record leaves the subsets, permutation conventions, and pointwise map unspecified, so a complete precise conjecture requires further source material. Its open status is supplied dated curation.
\subsection{Sources}
\begin{itemize}
\item Local UnsolvedMath registry, supplied record with permanent ID \(30004937\), title \emph{Shift Maps Between Totally Nonnegative Grassmannians}, exact source problem code \texttt{catalogue entry}. The supplied registry labels the problem \texttt{open} and its statement \texttt{corrected\_verified}; these are supplied labels, not independent verification by this draft.
\item Supplied report citation: \emph{Mini-Workshop: Scattering Amplitudes, Cluster Algebras, and Positive Geometries} (2022), Oberwolfach Reports, pp.\ 3101--3133, DOI \texttt{10.4171/owr/2021/57}. \url{https://doi.org/10.4171/owr/2021/57}. The supplied excerpt is attributed to Pavel Galashin under the contribution title \emph{Shift maps, poset associahedra, and totally nonnegative critical varieties}; it includes a page marker for p.\ 3107. The registry records proposed year 2021 and provides no separate proposer attribution.
\item Supplied literature source: ``Galashin, totally nonnegative critical varieties.'' \url{https://academic.oup.com/imrn/advance-article/doi/10.1093/imrn/rnad084/7146963}.
\item Supplied literature source: ``original Oberwolfach report.'' \url{https://doi.org/10.4171/owr/2021/57}.
\item Literature assessment dated exactly \texttt{2026-08-22}: the supplied curation describes the Grassmannian shift-map bijection as conjectural and states that later results for critical varieties do not establish the full stated Grassmannian bijection. This is supplied dated curation, not independently verified literature assessment. The supplied background explicitly requests independent verification of time-sensitive status and cited claims.
\item Unresolved internal report references: [4, 5, 6] concern critical varieties and associated polytopes; [11] concerns the name T-duality; [12] concerns the conjectural Grassmannian bijection; [2] and [1] concern the BCFW triangulation and amplituhedron; [9] and [7, Question 9.2] concern the proposed network-space topological equivalence. Their bibliographic entries are absent from the supplied record.
\end{itemize}
\end{document}
